\documentclass[11pt]{article}%设置文档类型{article}和字体大小{11pt}，适配本地字体

\usepackage[a4paper]{geometry}%页面布局，纸张大小[a4paper]
\geometry{left=2.0cm,right=2.0cm,top=2.5cm,bottom=2.5cm}%页面边距
\usepackage[UTF8]{ctex}%latex支持中文宏包
\usepackage{amsmath, amsfonts, amssymb, amsthm}%数学公式环境、数学字体、数学字体扩展包、定理环境
\usepackage{bm,mathtools}%数学粗体、amsmath扩展包
\usepackage{graphicx}%插入和管理图片
\usepackage[graphicx]{realboxes}%创建和管理盒子（装有文本和图形）
\usepackage{algorithm,algorithmicx}%支持伪代码宏包
\usepackage[noend]{algpseudocode}%支持伪代码宏包
\usepackage{fancyhdr}%定制页眉和页脚
\setlength{\headheight}{13.6pt}%避免页眉高度警告
\pagestyle{fancy}
\fancyhf{} % 清空默认的页眉和页脚
\fancyhead[L]{\kaishu 中国科学院大学}
\fancyhead[C]{林俊杰2023K8009916012}
\fancyhead[R]{\kaishu}
\fancyfoot[C]{\thepage}
\usepackage{booktabs}%美化表格宏包
\usepackage{caption}%定制图表标题宏包
\usepackage{xcolor}%切换字体颜色
\usepackage{placeins}%控制浮动体位置
\usepackage{cleveref}%智能引用
\crefname{figure}{图}{figures}
\Crefname{figure}{Figure}{Figures}
\crefname{table}{表}{tables}
\Crefname{table}{Table}{Tables}
\usepackage{tikz} % 绘图宏包
\usetikzlibrary{arrows.meta, positioning, calc, shapes, fit, decorations.pathreplacing} % 加载绘图库   

%数学环境计数器
\newcounter{counter_exm}\setcounter{counter_exm}{1}
%\newcounter{counter_thm}\setcounter{counter_thm}{1}
%\newcounter{counter_lma}\setcounter{counter_lma}{1}
%\newcounter{counter_dft}\setcounter{counter_dft}{1}
%\newcounter{counter_clm}\setcounter{counter_clm}{1}
%\newcounter{counter_cly}\setcounter{counter_cly}{1}

%图标随章节计数{章节.张数}
\counterwithin{figure}{section}
\counterwithin{table}{section}

\newtheorem{theorem}{{\hskip 1.7em \bf 定理}}
\newtheorem{lemma}[theorem]{\hskip 1.7em 引理}
\newtheorem{proposition}[theorem]{Proposition}
\newtheorem{claim}[theorem]{\hskip 1.7em 命题}
\newtheorem{corollary}[theorem]{\hskip 1.7em 推论}
\newtheorem{definition}[theorem]{\hskip 1.7em 定义}

\renewcommand{\emph}[1]{{\kaishu #1}}
\newcommand{\articlename}{生成函数的妙用}

\newenvironment{solution}{{\noindent\hskip 2em \bf 解 \quad}}{\par}
\newenvironment{example}{{\noindent\hskip 2em \bf 例 \arabic{counter_exm}\quad}}{\addtocounter{counter_exm}{1}\par}
\newenvironment{concept}[1]{{\bf #1\quad} \begin{kaishu}}{\end{kaishu}\par}
\renewenvironment{proof}{{\noindent\hskip 2em \bf 证明 \quad}}{\hfill$\qed$\par}

\begin{document}
    \begin{center}
        \LARGE \bf \articlename
    \end{center}

计算
\[
\lim_{n \to \infty} \int_0^1 (1 - x^2)^n \, dx
\]

\begin{solution}
令 
\[
a_n = \int_0^1 (1 - x^2)^n \, dx
\]

构造生成函数
\[
\begin{aligned}
F(s) &= \sum_{n=0}^\infty a_n \cdot s^n \quad (0 < s < 1) \\
&= \sum_{n=0}^\infty \int_0^1 (1 - x^2)^n \, dx \cdot s^n \\
&= \int_0^1 \sum_{n=0}^\infty \bigl((1 - x^2)^n s^n\bigr) \, dx \\
&= \int_0^1 \frac{dx}{1 - (1 - x^2) s} \\
&= \frac{\arcsin \sqrt{s}}{\sqrt{s(1 - s)}}
\end{aligned}
\]

注意到
\[
F'(s) = \frac{1 - \frac{(1 - 2s) \cdot \arcsin \sqrt{s}}{\sqrt{s(1 - s)}}}{2s(1 - s)} = \frac{1 - (1 - 2s) F(s)}{2s(1 - s)}
\]

整理得
\[
2s(1 - s) F'(s) + (1 - 2s) F(s) = 1
\]

这是一个关于 \( F(s) \) 的微分方程，将 \( F(s) \) 展开成幂级数
\[
2s(1 - s) \sum_{n=0}^\infty n \cdot a_n \cdot s^{n-1} + (1 - 2s) \sum_{n=0}^\infty a_n \cdot s^n = 1
\]

再整理得
\[
a_0 + (2a_1 + a_1 - 2a_0) \cdot s + \sum_{n=2}^\infty \bigl((2n+1)a_n - 2n \cdot a_{n-1}\bigr) \cdot s^n = 1
\]

对比系数得
\[
a_0 = 1,\quad a_1 = \frac{2}{3},\quad \frac{a_n}{a_{n-1}} = \frac{2n}{2n+1}
\]

因此
\[
a_n = \frac{(2n)!!}{(2n+1)!!}
\]

接下来的放缩利用
\[ 
n > \sqrt{n^2 - 1} = \sqrt{n - 1} \cdot \sqrt{n + 1} 
\]

因此
\[
\begin{aligned}
a_n &= \frac{2 \times 4 \times 6 \times 8 \times \cdots \times 2n}{3 \times 5 \times 7 \times 9 \times \cdots \times (2n+1)} \\
&\le \frac{2 \times 4 \times 6 \times 8 \times \cdots \times 2n}{\sqrt{2} \times \sqrt{4} \times \sqrt{4} \times \sqrt{6} \times \sqrt{6} \times \sqrt{8} \times \sqrt{8} \times \sqrt{10} \times \cdots \times \sqrt{2n} \cdot \sqrt{2n+2}} \\
&= \frac{2}{\sqrt{2} \cdot \sqrt{2n+2}} \\
&= \frac{1}{\sqrt{n+1}}
\end{aligned}
\]

因为
\[
\lim_{n \to \infty} \frac{1}{\sqrt{n+1}} = 0
\]

所以
\[
\lim_{n \to \infty} \int_0^1 (1-x^2)^n \, dx = \lim_{n \to \infty} a_n = 0
\]
\end{solution}
\end{document}